\documentclass[aspectratio=169]{beamer}

% Build with: tectonic talk/slides.tex

\usetheme{default}
\usefonttheme{professionalfonts}
\setbeamertemplate{navigation symbols}{}
\setbeamertemplate{footline}[frame number]
\setbeamertemplate{itemize items}[circle]

\usepackage{listings}
\usepackage{xcolor}

\definecolor{accent}{RGB}{183,65,14}      % rust orange
\definecolor{keyword}{RGB}{0,92,153}
\definecolor{comment}{RGB}{96,128,96}
\definecolor{string}{RGB}{136,0,0}
\definecolor{macro}{RGB}{122,62,157}
\definecolor{codebg}{RGB}{247,247,245}

\setbeamercolor{frametitle}{fg=accent}
\setbeamercolor{title}{fg=accent}
\setbeamercolor{structure}{fg=accent}
\setbeamerfont{frametitle}{series=\bfseries}
\setbeamertemplate{blocks}[rounded][shadow=false]
\setbeamercolor{block title alerted}{fg=string,bg=string!12}
\setbeamercolor{block body alerted}{fg=black,bg=string!4}

\lstdefinelanguage{Rust}{
  keywords={as,break,const,continue,crate,dyn,else,enum,extern,false,fn,for,
    if,impl,in,let,loop,match,mod,move,mut,pub,ref,return,static,struct,
    super,trait,true,type,unsafe,use,where,while,async,await,self,Self},
  sensitive=true,
  comment=[l]{//},
  string=[b]",
  alsoletter={\#},
}

\lstset{
  language=Rust,
  basicstyle=\ttfamily\footnotesize,
  keywordstyle=\color{keyword}\bfseries,
  commentstyle=\color{comment}\itshape,
  stringstyle=\color{string},
  backgroundcolor=\color{codebg},
  frame=single,
  framerule=0pt,
  xleftmargin=1em,
  framexleftmargin=1em,
  showstringspaces=false,
  columns=fullflexible,
  keepspaces=true,
  moredelim=[s][\color{macro}]{\#[}{]},
}

\title{\texttt{dyn\_shim}}
\subtitle{Trait objects for traits that refuse}
\author{Nathan Lilienthal}
\date{\small\texttt{crates.io/crates/dyn\_shim}}

\begin{document}

\begin{frame}
  \titlepage
\end{frame}

% ---------------------------------------------------------------- 1. problem

\begin{frame}[fragile]{The goal: a mixed bag of implementors}
\begin{lstlisting}
trait Channel {
    fn label(&self) -> String;
    fn deliver(&mut self, message: &str);
    fn connect() -> Self;
}

let channels: Vec<Box<dyn Channel>> = vec![
    Box::new(Email::connect()),
    Box::new(Webhook::connect()),
];
\end{lstlisting}
\pause
\begin{lstlisting}[language={},basicstyle=\ttfamily\footnotesize\color{string}]
error[E0038]: the trait `Channel` is not dyn compatible
  = note: because associated function `connect` has no `self` parameter
\end{lstlisting}
\end{frame}

\begin{frame}{Why: a trait object is a vtable}
  \begin{itemize}
    \item \texttt{Box<dyn Trait>} is a data pointer plus a table of function
          pointers, one per method.
    \item Every method has to fit in one vtable entry. These cannot:
    \begin{itemize}
      \item no receiver (\texttt{fn connect() -> Self})
      \item \texttt{Self} by value in arguments or returns
      \item generic methods (\texttt{fn explain<W: Write>})
      \item \texttt{impl Trait} in the signature
      \item supertraits like \texttt{Clone} (whose \texttt{clone} returns \texttt{Self})
    \end{itemize}
    \item And one bad method poisons the \emph{whole trait}, not just itself.
  \end{itemize}
\end{frame}

% ------------------------------------------------- 2. the standard escape hatch

\begin{frame}[fragile]{The standard escape hatch: \texttt{where Self: Sized}}
\begin{lstlisting}
trait Channel {
    fn label(&self) -> String;
    fn deliver(&mut self, message: &str);
    fn connect() -> Self where Self: Sized;
}

let channels: Vec<Box<dyn Channel>> = /* works! */;
\end{lstlisting}
  \begin{itemize}
    \item The gated method is hidden from \texttt{dyn Channel}, so the trait is
          dyn-compatible again.
    \item Still callable on every concrete type.
    \item Honestly: this covers a lot of the basic cases.
  \end{itemize}
\end{frame}

\begin{frame}{Where \texttt{Self: Sized} runs out}
  \begin{enumerate}
    \item<1-> \textbf{It is invasive.} dyn-compatibility bookkeeping lives on
          your source trait's API forever.
    \item<2-> \textbf{The method could work on an erased value, with a tweak.}
          \texttt{fn close(self)} could ride as \texttt{Box<Self>};
          \texttt{fn explain<W: Write>} could take \texttt{\&mut dyn Write}.
          Gating just throws them away.
    \item<3-> \textbf{The problem is a supertrait.}
          \texttt{trait Event: Clone} has no method to gate.
    \item<4-> \textbf{It is not your trait.} You cannot edit a dependency's
          trait to add the gate.
  \end{enumerate}
\end{frame}

% ----------------------------------------------------------------- 3. the core

\begin{frame}[fragile]{The core idea: generate a second trait}
  You write:
\begin{lstlisting}
#[dyn_shim(DynChannel)]
trait Channel {
    fn label(&self) -> String;
    // ...
}
\end{lstlisting}
\pause
  The macro generates:
\begin{lstlisting}
trait DynChannel {
    fn label(&self) -> String;
}

impl<T: Channel> DynChannel for T {
    fn label(&self) -> String { Channel::label(self) }
}
\end{lstlisting}
  \texttt{Channel} stays exactly as written, and the blanket impl means every
  \texttt{Channel} is already a \texttt{DynChannel}.
\end{frame}

\begin{frame}[fragile]{No receiver: skipped}
\begin{lstlisting}
#[dyn_shim(DynChannel)]
trait Channel {
    // ...
    fn connect() -> Self;   // receiverless
}
\end{lstlisting}
\pause
\begin{lstlisting}
trait DynChannel {
    // ...
    // no connect: it cannot ride a vtable
}
\end{lstlisting}
  \begin{itemize}
    \item What \texttt{Self: Sized} did by gating, the shim does by omission,
          and the source trait needs no edit.
    \item \texttt{connect} is still right there on every concrete type.
  \end{itemize}
\end{frame}

\begin{frame}[fragile]{By-value \texttt{self}: rewritten}
\begin{lstlisting}
#[dyn_shim(DynChannel)]
trait Channel {
    // ...
    fn close(self) -> u32;  // by-value self
}
\end{lstlisting}
\pause
\begin{lstlisting}
trait DynChannel {
    // ...
    fn close(self: Box<Self>) -> u32;   // self -> Box<Self>
}

impl<T: Channel> DynChannel for T {
    // ...
    fn close(self: Box<Self>) -> u32 { Channel::close(*self) }
}
\end{lstlisting}
  Not skipped: \emph{rewritten} into a form a vtable can carry.
\end{frame}

\begin{frame}[fragile]{A single generic argument: erased}
\begin{lstlisting}
#[dyn_shim(DynChannel)]
trait Channel {
    // ...
    #[dyn_shim(erase)]
    fn dump<W: Write>(&self, out: &mut W);  // generic argument
}
\end{lstlisting}
\pause
\begin{lstlisting}
trait DynChannel {
    fn dump(&self, out: &mut dyn Write);    // W erased
}
impl<T: Channel> DynChannel for T {
    fn dump(&self, mut out: &mut dyn Write) {
        Channel::dump(self, &mut out)       // W = &mut dyn Write
    }
}
\end{lstlisting}
  \begin{itemize}
    \item Generic behind \texttt{\&}/\texttt{\&mut}, single bound: forwards
          for real.
    \item Opt-in: \texttt{W} now dispatches through \texttt{dyn Write} inside
          the impl.
  \end{itemize}
\end{frame}

\begin{frame}[fragile]{A method returning \texttt{Self}: boxed}
\begin{lstlisting}
#[dyn_shim(DynChannel)]
trait Channel {
    // ...
    #[dyn_shim(boxed)]
    fn muted(&self) -> Self;  // Self return
}
\end{lstlisting}
\pause
\begin{lstlisting}
trait DynChannel {
    fn muted(&self) -> Box<dyn DynChannel>;  // Self boxed
}
impl<T: Channel> DynChannel for T {
    fn muted(&self) -> Box<dyn DynChannel> {
        Box::new(Channel::muted(self))
    }
}
\end{lstlisting}
  \begin{itemize}
    \item Opt-in: the one rewrite that adds a cost, one allocation per call.
    \item Bare \texttt{-> Self} only; \texttt{Option<Self>} and friends stay
          skipped for now.
  \end{itemize}
\end{frame}

\begin{frame}[fragile]{Using it}
\begin{lstlisting}
let channels: Vec<Box<dyn DynChannel>> = vec![
    Box::new(Email::connect()),    // concrete side: connect works
    Box::new(Webhook::connect()),
];

for ch in &channels {
    println!("{}", ch.label());    // forwarded
}

for ch in channels {
    let sent = ch.close();         // by-value, through Box<Self>
}
\end{lstlisting}
  \begin{itemize}
    \item Even the by-value \texttt{close(self)} works on the erased object,
          which \texttt{Self: Sized} could never give you.
  \end{itemize}
\end{frame}

% -------------------------------------------------------------- 4. reflexive

\begin{frame}[fragile]{The shim's debt: \texttt{reflexive}}
  Gating kept identity free: \texttt{dyn Channel} is a \texttt{Channel}.
  A shim is a \emph{distinct} trait; source-generic code rejects its objects:
\begin{lstlisting}
// ?Sized: lets R be dyn DynRule
fn passes<R: Rule + ?Sized>(rule: &R, v: i32) -> bool { ... }
\end{lstlisting}
\pause
  \texttt{reflexive} buys it back, implementing \texttt{Rule} for the shim's
  objects:
\begin{lstlisting}
// bare: impl Rule for dyn DynRule; boxed: for Box<dyn DynRule>
#[dyn_shim(DynRule, reflexive = bare + boxed)]
trait Rule {
    fn check(&self, value: i32) -> bool;
    #[dyn_shim(stub = None)]        // degrade gracefully
    fn parsed<T: FromStr>(&self, text: &str) -> Option<T>;
}
\end{lstlisting}
\pause
  An erased value flows back in:
\begin{lstlisting}
let rule: Box<dyn DynRule> = ...;
passes(&*rule, 12);
\end{lstlisting}
\end{frame}

\begin{frame}{Reflexive: accounting for every method}
  The reflexive impl must account for \emph{every} source method:
  \begin{itemize}
    \item Whatever the shim carries forwards through it, including the
          \texttt{erase} and \texttt{boxed} rewrites from before: nothing new.
    \item \texttt{\#[dyn\_shim(stub = expr)]}: cannot forward, so return a
          fallback value instead.
    \item \texttt{\#[dyn\_shim(panic)]}: the last resort, when no fallback
          makes sense. Just sugar for \texttt{stub = panic!("...")} with a
          message naming the method and shim.
  \end{itemize}
  \begin{alertblock}{The honest tradeoff}
    \texttt{where Self: Sized} rejects a bad call at \emph{compile time};
    \texttt{\#[dyn\_shim(panic)]} compiles and aborts at \emph{runtime}.
    If hiding the method is the only remediation a trait needs, gating
    fails earlier. Reach for \texttt{panic} only when the reflexive impl
    is worth that risk.
  \end{alertblock}
\end{frame}

% ------------------------------------------------------------ 5. clone / hash

\begin{frame}[fragile]{A \texttt{Clone} supertrait: nothing to gate}
\begin{lstlisting}
trait Shape: Clone {  // not dyn compatible: clone returns Self
    fn area(&self) -> u32;
}
\end{lstlisting}
  \begin{itemize}
    \item Nothing to gate: the offender is a \emph{bound}, not a method.
    \item The classic fix is a whole crate per trait:
          \texttt{dyn-clone} with \texttt{clone\_trait\_object!},
          \texttt{dyn-hash} with \texttt{hash\_trait\_object!}.
  \end{itemize}
\end{frame}

\begin{frame}[fragile]{Recognized bounds}
\begin{lstlisting}
#[dyn_shim(DynShape: Clone + Hash)]
trait Shape {
    fn area(&self) -> u32;
    fn scale(&mut self, factor: u32);
}

let shapes: Vec<Box<dyn DynShape>> = /* ... */;
let copy = shapes.clone();        // Box<dyn DynShape>: Clone
let hash = RandomState::new().hash_one(&*shapes[0]); // dyn DynShape: Hash
\end{lstlisting}
  \begin{itemize}
    \item \texttt{Clone} and \texttt{Hash} in the shim's bounds are
          \emph{recognized}: instead of becoming supertraits, they generate
          cloning and hashing machinery on the trait objects.
    \item The same move as \texttt{reflexive}, aimed at a \emph{bound}:
          implement the unreachable trait on the shim's trait objects.
    \item No crate features, no \texttt{unsafe}, written or generated.
  \end{itemize}
\end{frame}

\begin{frame}[fragile]{The \texttt{dyn-clone} / \texttt{dyn-hash} replacement}
  Already have a dyn-compatible trait? Bind the shipped
  \texttt{DynClone} / \texttt{DynHash} onto it:
\begin{lstlisting}
use dyn_shim::{dyn_shim_bind, DynClone, DynHash};

#[dyn_shim_bind(DynClone + DynHash)]
trait Product: DynClone + DynHash {
    fn price(&self) -> u32;
}
\end{lstlisting}
  \begin{itemize}
    \item Near drop-in for \texttt{dyn-clone} + \texttt{dyn-hash}: keep the
          supertraits, replace the \texttt{*\_trait\_object!} calls with one
          attribute.
    \item The point: those crates are each a hand-built special case.
          Here they fall out of the general mechanism, and any
          \texttt{\#[dyn\_shim]} shim binds the same way.
    \item The crate's one \texttt{unsafe} block lives here: cloning through
          an already-erased object needs a pointer-metadata splice, taken
          from \texttt{dyn-clone} itself.
  \end{itemize}
\end{frame}

% -------------------------------------------------------------- 6. foreign

\begin{frame}[fragile]{Not your trait? \texttt{\#[dyn\_shim\_foreign]}}
\begin{lstlisting}
use dyn_shim::dyn_shim_foreign;

// vendor::Widget has build() -> Self and consume(self).
#[dyn_shim_foreign(vendor::Widget)]
trait DynWidget: Clone {  // recognized, like DynShape's Clone
    fn render(&self) -> String;
    fn resize(&mut self, by: i32);
    fn consume(self) -> String;   // same Box<Self> rewrite
}
\end{lstlisting}
  \begin{itemize}
    \item A proc macro cannot see another crate's source, so here the
          annotated trait \emph{is} the shim: restate what should forward.
    \item A restated signature that does not match fails to compile.
  \end{itemize}
\end{frame}

% -------------------------------------------------------------------- 7. close

\begin{frame}{Recap: one mechanism, three paths}
  \begin{itemize}
    \item \texttt{where Self: Sized}: when hiding the offending methods is
          enough, stop there.
    \item \texttt{\#[dyn\_shim]}: a generated shim trait; skip or rewrite
          what a vtable cannot carry, keep the source trait clean.
  \end{itemize}
  \vspace{0.5em}
  Three paths from there, in any combination:
  \begin{itemize}
    \item \texttt{reflexive}: shim objects satisfy the source trait again
    \item recognized \texttt{Clone} / \texttt{Hash} bounds and
          \texttt{\#[dyn\_shim\_bind]}: the \texttt{dyn-clone} /
          \texttt{dyn-hash} problem, generalized
    \item \texttt{\#[dyn\_shim\_foreign]}: traits you do not own
  \end{itemize}
\end{frame}

\begin{frame}{Thanks}
  \begin{center}
    {\Large \texttt{dyn\_shim}}\\[1em]
    \texttt{crates.io/crates/dyn\_shim}\\
    \texttt{docs.rs/dyn\_shim}\\
    \texttt{github.com/nixpulvis/dyn\_shim}\\[1em]
    One runnable example per feature in \texttt{examples/}:\\
    \texttt{shim}, \texttt{clone\_and\_hash}, \texttt{bind},
    \texttt{reflexive}, \texttt{foreign}
  \end{center}
\end{frame}

\end{document}
